\documentclass[10pt,a4paper]{amsart}
\usepackage[margin=19mm]{geometry}
\usepackage{amsmath,amssymb,mathtools}
\usepackage[scr=rsfs,cal=euler]{mathalfa}
\usepackage{xfrac}
\usepackage[unicode,hidelinks,bookmarks=false]{hyperref}
\newcommand{\ie}{i.e.\ }
\newcommand{\cf}{cf.\ }
\newcommand{\resp}{respectively\ }
\newcommand{\Subsection}{\subsection}
\providecommand{\nobreakdash}{\nobreak\hbox{-}}
\setlength{\emergencystretch}{2em}
\multlinegap0pt
\usepackage{amssymb}
\usepackage[shortlabels]{enumitem}
\usepackage{csquotes}
\usepackage{xcolor}
\newtheorem{theorem}{Theorem}
\title{Generalized Kolmogorov systems with applications to astrophysics and biology}
\author{Dorota Bors, Robert Stańczy}
\date{Source: arXiv:2604.09720v1 (CC BY 4.0)}
\begin{document}
\maketitle

\noindent\emph{Source and licence.} Generalized Kolmogorov systems with applications to astrophysics and biology. Version arXiv:2604.09720v1 (CC BY 4.0), distributed by arXiv under the Creative Commons Attribution 4.0 International licence, http://creativecommons.org/licenses/by/4.0/, as recorded in the arXiv licence metadata for this exact version and shown on the abstract page https://arxiv.org/abs/2604.09720v1. Full licence notice: \texttt{licenses/arxiv-2604.09720-CC-BY-4.0.txt}.\par\smallskip

\noindent\textit{Editorial scope.} Counted unit: the article's theorem MyThmLabel (source0.tex lines 75--79). The article's complete original proof of it is delivered with it (source0.tex lines 80--90, one \texttt{proof} environment of 11 lines). No mathematics is written, repaired or re-derived: the statement and the proof are the article's own lines, in the article's own notation, and no retained line is substituted or re-flowed.  Retained uncounted as the source-local context the proof needs: lines 65--70.  Nothing was omitted from any retained span, so the package carries no omission record.  No retained line contains a citation, so the wrapper carries no bibliography and no bibliography entry.  No retained line contains a figure environment, an image-inclusion command or drawing code, so no image is delivered and the package declares no asset dependency.  Editorial formatting only, in addition: an AMS \texttt{amsart} preamble that restates the article's own declarations and macros used by the retained lines, namely the environments \texttt{theorem} on separate counters, together with the standard packages those lines need.  The retained environments print in this document as Theorem 1 in order of appearance, whereas the article prints the same items with the numbering of its own full document.  The complete native source of the article as distributed in the arXiv e-print for this exact version is bundled unchanged as \texttt{sources/198225/source0.tex}.
\begin{equation}\label{Kol}
	\begin{cases}
		x'=g(x)G(x,y),\\
		y'=h(y)H(x,y)\,.
	\end{cases}
\end{equation}

\begin{theorem} \label{MyThmLabel} %Please insert your own label's name here.
Assume that $x,y \ge 0$. Suppose that $G_x\cdot H_x^z \le 0$ and $H_y\cdot G_y^w\ge 0$, where we denote  $H^z=H(x,z)$ and $G^w=G(w,y)$, while with subscripts partial derivatives are denoted. Then, the function
\[L(x,y)={\mathcal H}(x)+{\mathcal G}(y)\] 
is a Lyapunov function for (\ref{Kol}), where ${\mathcal H}(w)={\mathcal G}(z)=0$ and ${\mathcal G, H}>0$ elsewhere. The derivatives of these functions are defined as ${\mathcal H}'(x)=H(x,z)/g(x)$, ${\mathcal G}'(y)=-G(w,y)/h(y)$. 
\end{theorem}

\begin{proof}
Multiplying the first equation of the system (\ref{Kol}) by $H(x,z)/g(x)$ and the second one by $-G(w,y)/h(y)$, and summing the results, yields
\[
x'H(x,z)/g(x)-y'G(w,y)/h(y)=G(x,y)(H^z-H(w,z))-H(x,y)(G^w-G(w,z))\,.
\]
By adding and subtracting the term $G(w,y)H(x,z)$ on the right hand side, and using the Mean Value Theorem alongside the definitions of ${\mathcal G}$ and ${\mathcal H}$, the derivative of the Lyapunov function along the orbit of (\ref{Kol}) $(L(x,y))'=x'L_x+y'L_y=x'{\mathcal G}'(x)+{\mathcal H}'(y)$, can be rephrased as
\[
(L(x,y))'=G_xH_x^z(x-w)^2-H_yG_y^w(y-z)^2\,,
\]
Given our initial assumptions regarding the signs of these partial derivatives, it follows that $(L(x,y))'\le 0$, which completes the proof.
\end{proof}

\end{document}
